\subsection{Hilbert Space and Operator Contents}

In a General CFT, the Hilbert Space and all Operator Contents can be
classified into the Primaries and Their Descendent.
\imp{Hilbert Space}{
  The Hilbert Space $ \mathcal{H} $ of a CFT is spanned by the
  highest weight representations of the Virasoro Algebra. Each
  highest weight representation is generated by a Primary State $
  |h,\bar{h}\rangle $ and its Descendent States $
  L_{-n_1}L_{-n_2}...|h,\bar{h}\rangle  $ ($ n_i>0 $).
  \begin{align}
    \mathcal{H} = \bigoplus_{h,\bar{h}} \mathcal{V}_h \otimes
    \bar{\mathcal{V}}_{\bar{h}} ,
  \end{align}
}

And the Operators are classified similarly
\imp{Operator Content}{
  The Operator Content of a CFT is classified by the Primary
  Operators $ \phi_{h,\bar{h}}(z,\bar{z}) $ and their Descendent
  Operators $ \phi_{h,\bar{h}}^{(-k_1,...,-k_n)}(z,\bar{z})  $ ($ k_i>0 $).
}

\subsection{Correlation Functions}\label{sec:CFT Correlation Functions}

To know all correlation functions of a CFT is equivalent to knowing
all physical observables of the CFT which is the ultimate goal of
studying a CFT. In fact the coorelation function of a CFT is fully
determined by the OPEs which we also called the \textbf{Operator
Algebra} of the CFT.
\defi{
  Operator Algebra

  The Operator Algebra is the collection of all OPEs between the
  operators in the CFT.
}
If we know the Operator Algebra of a CFT, we can in principle reduce
any N-point Correlation Functions to 2-Point Correlation Functions
(which is fully determined by conformal symmetry). Thus the full
information of a CFT correlation function is encoded in its Operator Algebra.

In fact, the operator algebra of a CFT can be determined if we know a
series of data, we write this as the following theorem

\thm{
  Determination of Operator Algebra

  The Operator Algebra of a CFT is fully determined if we know the
  following data:
  \begin{itemize}
    \item Central Charge $ c $
    \item Hilbert Space or the Operator Content
    \item Structural Constants $ C_{abc} $
  \end{itemize}
}
Proof: See section 6.6.3 of the yellow book for reference. It just
calculate the OPE and find that it only depends on these data.

Then the question is, which combination of these data can give a
self-consistent CFT? This is the core idea of Conformal Bootstrap
which we will discuss in the next section.

\subsection{Axiomatic Formalism of CFT}

The above three sections we explicitly and implicitly made some
assumptions apart from the assumption of conformal symmetry :
\begin{itemize}
  \item A CFT should contain Primary Fields and their descendants only
  \item General CFT should contain a E-M Tensor with a suitable OPE
  \item Conformal Ward identity holds beyond primary fields
  \item A global conformal invariant vacuum $ |0\rangle $
  \item Hilbert Space is spanned by the highest weight
    representations of the Virasoro Algebra
  \item ...
\end{itemize}
These assumptions are really messy. However, they made the final
structure of CFT beautiful and clear. Thus, we generally won't use
these explicit and implicit assumptions in the above sections to
define a CFT (or we will fall into mess again), but we will make the
clear results (OPE Structure, Virasoro Algebra Representation
Spectrum , state operator correspondence...) as a defining feature of
a CFT. These feature gives a abstract but clear definition of a CFT,
this work is first done by Belavin, Polyakov and Zamolodchikov in
their famous BPZ paper and refined through modern works.
